\documentclass[../readings.tex]{subfiles}

\begin{document}

\subsection{Fourier Analysis as Change of Basis}
\label{sub:fourier}

\textBox{%
Fourier analysis provides a framework for representing functions and data in terms of frequency components. In engineering, this transformation is critical for signal processing, image compression, and solving differential equations where the Laplacian operator is diagonalized by the Fourier basis.
}

\subsubsection{Discrete Fourier Transform (DFT)}

\addDef{Periodic Grid Signal}{%
A discrete signal $\bx\in\fC^n$ is interpreted as one period of a periodic sequence:
\begin{align*}
    x_{j+n}=x_j.
\end{align*}
The DFT represents this signal in the basis of discrete complex exponentials.
}
\label{def:periodic-grid-signal}

\addDef{DFT Matrix}{%
Unitary matrix $\mathbf{F} \in \fC^{n \times n}$ with entries:
\begin{align*}
    F_{jk} = \frac{1}{\sqrt{n}} \omega_n^{jk}, \quad \omega_n = e^{-2\pi i/n}.
\end{align*}
\begin{itemize}
    \item \textbf{Unitary Property:}\enspace $\mathbf{F}^* \mathbf{F} = \bI$. Orthonormal columns.
    \item \textbf{Inverse DFT:}\enspace $\mathbf{F}^{-1} = \mathbf{F}^*$.
\end{itemize}
}
\label{def:dft-matrix}

\addThm{Parseval's Theorem}{%
The DFT is an isometry: $\|\mathbf{F}\bx\|_2 = \|\bx\|_2$. Total energy is preserved between time and frequency domains.
}
\label{thm:dft-parseval}

\addPrf{%
Since $\mathbf{F}$ is unitary, $\mathbf{F}^*\mathbf{F}=\bI$. Hence
\begin{align*}
    \|\mathbf{F}\bx\|_2^2
    =
    (\mathbf{F}\bx)^*(\mathbf{F}\bx)
    =
    \bx^*\mathbf{F}^*\mathbf{F}\bx
    =
    \bx^*\bx
    =
    \|\bx\|_2^2.
\end{align*}
Taking square roots gives the claim.
}

\addExcr{%
\begin{enumerate}
    \item Use Def~\ref{def:dft-matrix} to compute the $(j,\ell)$ entry of $\mathbf{F}^*\mathbf{F}$.
    \item Evaluate the finite geometric sum that appears.
    \item Conclude $\mathbf{F}^*\mathbf{F}=\bI$.
    \item Explain why this is the finite Fourier analog of orthonormal coordinates.
\end{enumerate}
}
\label{ex:dft-parseval}

\addDef{Frequency Bins}{%
For samples on $[0,2\pi)$, the integer frequency attached to index $k$ is usually ordered as
\begin{align*}
    0,1,2,\dots,\lfloor n/2\rfloor,-\lfloor(n-1)/2\rfloor,\dots,-1.
\end{align*}
This is the ordering returned by \texttt{np.fft.fftfreq} after scaling by $n$.
}
\label{def:frequency-bins}

\addRemark{%
\textbf{(Normalization convention)}\enspace Some libraries put the factor $1/n$ in the inverse transform instead of using the unitary factor $1/\sqrt n$ in both directions. The formulas change by constants, but the frequency content is the same.
}

\addDef{Nyquist Frequency and Aliasing}{%
On $n$ equally spaced samples of a $2\pi$-periodic signal, frequencies that differ by multiples of $n$ are indistinguishable:
\begin{align*}
    e^{i(k+n)x_j}=e^{ikx_j}.
\end{align*}
The largest resolvable unsigned frequency is the Nyquist frequency, approximately $n/2$ modes per period.
}
\label{def:nyquist-aliasing}

\addExcr{%
\begin{enumerate}
    \item Let $x_j=2\pi j/n$. Prove $e^{i(k+n)x_j}=e^{ikx_j}$ for every grid point.
    \item For $n=8$, list the frequency bin ordering from Def~\ref{def:frequency-bins}.
    \item Sample $\sin(7x)$ on $n=8$ points and identify its aliased frequency.
\end{enumerate}
}
\label{ex:nyquist-aliasing}

\subsubsection{Convolution and Circulant Matrices}

\addDef{Circulant Matrix}{%
Matrix $\mathbf{C}$ where each row is a cyclic shift of the first column $\bc$.
\begin{itemize}
    \item \textbf{Diagonalization:}\enspace Every circulant matrix is diagonalized by the DFT: $\mathbf{C} = \mathbf{F}^* \text{diag}(\mathbf{F}\bc) \mathbf{F}$.
    \item \textbf{Eigenvalues:}\enspace The eigenvalues of $\mathbf{C}$ are the DFT coefficients of its first column.
\end{itemize}
}
\label{def:circulant-matrix}

\addThm{Convolution Theorem}{%
Circular convolution $\bz = \bc * \bx$ is equivalent to pointwise multiplication in frequency:
\begin{align*}
    \mathbf{F}(\bc * \bx) = (\mathbf{F}\bc) \odot (\mathbf{F}\bx).
\end{align*}
\textbf{Complexity:}\enspace Reduced from $O(n^2)$ (direct) to $O(n \log n)$ (FFT).
}
\label{thm:convolution-theorem}

\addExmpl{%
\textbf{(Circular convolution for $n=3$)}\enspace If $\bc=(c_0,c_1,c_2)^T$, then
\begin{align*}
    \bc * \bx =
    \begin{pmatrix}
        c_0&c_2&c_1\\
        c_1&c_0&c_2\\
        c_2&c_1&c_0
    \end{pmatrix}
    \begin{pmatrix}x_0\\x_1\\x_2\end{pmatrix}.
\end{align*}
The wraparound terms are a consequence of the periodic convention in Def~\ref{def:periodic-grid-signal}.
}

\subsubsection{The Fast Fourier Transform (FFT)}

\addThm{Cooley-Tukey Algorithm}{%
Computes the DFT in \textbf{$O(n \log n)$} flops for $n = 2^p$.
\begin{itemize}
    \item \textbf{Mechanism:}\enspace Recursively splits $n$-point DFT into two $n/2$-point DFTs (even/odd indices).
    \item \textbf{NumPy:}\enspace Use \texttt{np.fft.fft} and \texttt{np.fft.ifft}.
\end{itemize}
}
\label{thm:cooley-tukey}

\addExcr{%
\begin{enumerate}
    \item Split the DFT sum into even and odd input indices.
    \item Show that each part is an $n/2$-point DFT multiplied by phase factors.
    \item Write the recurrence $T(n)=2T(n/2)+O(n)$.
    \item Solve the recurrence to obtain the cost in Thm~\ref{thm:cooley-tukey}.
\end{enumerate}
}
\label{ex:cooley-tukey}

\addRemark{%
\textbf{Spectral Differentiation:}\enspace For smooth periodic signals, differentiate by multiplying frequency components $\hat{x}_k$ by $ik$. Provides exponential accuracy compared to polynomial accuracy of finite differences.
}

\addThm{Fourier Differentiation Rule}{%
For a smooth $2\pi$-periodic function with Fourier series
\begin{align*}
    f(x)=\sum_k \hat f_k e^{ikx},
\end{align*}
the derivative has coefficients
\begin{align*}
    \widehat{f'}_k=ik\hat f_k.
\end{align*}
}
\label{thm:fourier-differentiation}

\subsubsection{Exercises}

\addExcr{%
\begin{enumerate}
    \item Write out the $4 \times 4$ DFT matrix explicitly. Verify it is unitary using Ex~\ref{ex:dft-parseval}.
    \item Solve the periodic Poisson equation $-u'' = f$ via FFT. (\textit{Hint: Laplacian is circulant}).
    \item Compare the speed of \texttt{np.fft.fft(x)} vs.\ \texttt{F @ x} for $n=4096$.
    \item Implement spectral differentiation for $f(x) = \sin(x)$ using Thm~\ref{thm:fourier-differentiation} and check error vs.\ $n$.
\end{enumerate}
}

\end{document}
